{\bf Prefaci}

\vspace{0.5em}
\noindent
Aquest document recull el treball de fi de grau realitzat durant el curs 2025--2026 en el marc del Grau en Enginyeria Matemàtica en Ciència de Dades de la Universitat Pompeu Fabra, en col·laboració amb el Dimension Lab de l'Hospital de la Santa Creu i Sant Pau.

\vspace{0.5em}
\noindent
El cos principal del treball (Introducció a Conclusions) està redactat en anglès per coherència amb la literatura científica del domini. Les seccions preliminars obligatòries (resums i agraïments) es presenten en català, anglès i castellà segons els requisits de la plantilla UPF.

\vspace{0.5em}
\noindent
\textbf{Declaració sobre l'ús de la intel·ligència artificial.}
D'acord amb la instrucció universitària sobre l'ús de la IA en activitats avaluables, declaro que el contingut tècnic i clínic del treball --- especialment la Introducció, la Metodologia i les Conclusions --- ha estat redactat i revisat per mi, amb el suport del meu director i supervisor. Per a la integració en \LaTeX{}, la correcció de format i l'esborrany de les seccions preliminars (agraïments, resums i prefaci), he utilitzat eines d'IA generativa com a assistent d'edició i maquetació. Tots els fragments generats o reformulats amb IA han estat revisats, validats i adaptats per mi abans de la seva inclusió. Qualsevol text, figura o referència inclosa sense revisió pròpia serà clarament identificada en versions futures del document si així ho requereix el director del treball.

\vspace{0.5em}
\noindent
\textbf{Nota sobre figures i dades.}
Les figures de fluxos clínics adaptades de treballs interns de l'hospital es citen explícitament al peu de figura. Les dades de cohorts clíniques i els resultats del pipeline es documentaran als capítols corresponents i, quan calgui, a l'apèndix de suport.
